% ============================================================================
% Bang dien giai tu viet tat va ky hieu — RIENG cho BC2 (do an tot nghiep).
% Tach khoi ../shared/abbreviations.tex (dung chung voi BC1) de bang ky hieu
% phan anh dung cac ky hieu dung trong bao cao nay, khong phai cua BC1.
% ============================================================================

\chapter*{BẢNG DIỄN GIẢI TỪ VIẾT TẮT VÀ TIẾNG ANH}
\addcontentsline{toc}{chapter}{BẢNG DIỄN GIẢI TỪ VIẾT TẮT VÀ TIẾNG ANH}

\begin{longtable}{@{}p{3.0cm} >{\raggedright\arraybackslash}p{6cm} >{\raggedright\arraybackslash}p{5.5cm}@{}}

\textbf{Tên viết tắt} & \textbf{Tên (English)} & \textbf{Diễn giải (Tiếng Việt)} \\
\midrule
\endfirsthead

BFS     & Breadth-First Search & Tìm kiếm theo chiều rộng \\
BPMN    & Business Process Model and Notation & Ký pháp mô hình hóa quy trình nghiệp vụ \\
BFD     & Business Function Diagram & Sơ đồ phân rã chức năng nghiệp vụ \\
BTF     & BPF Type Format & Định dạng thông tin kiểu dữ liệu của kernel, nền tảng cho CO-RE \\
C4      & C4 model & Mô hình kiến trúc phần mềm bốn mức (Context/Container/Component/Code) \\
CI      & Confidence Interval & Khoảng tin cậy (thống kê), ở đây dùng mức 95\% \\
CMS     & Count-Min Sketch & Cấu trúc dữ liệu "sketch" ước lượng tần suất phần tử trong luồng với sai số kiểm soát được \\
CNI     & Container Network Interface & Chuẩn cắm mạng cho container/Kubernetes (vd Calico, Cilium) \\
CO-RE   & Compile Once -- Run Everywhere & Kỹ thuật biên dịch eBPF một lần chạy trên nhiều phiên bản kernel nhờ BTF \\
CSDL    & Database & Cơ sở dữ liệu \\
DFD     & Data Flow Diagram & Sơ đồ luồng dữ liệu \\
eBPF    & extended Berkeley Packet Filter & Kỹ thuật chạy mã trong nhân Linux để quan sát/xử lý sự kiện một cách an toàn \\
ERD     & Entity--Relationship Diagram & Sơ đồ thực thể -- quan hệ (thiết kế CSDL) \\
FPR     & False Positive Rate & Tỷ lệ báo động giả (lưu lượng hợp lệ bị gắn cờ tấn công) \\
FR      & Functional Requirement & Yêu cầu chức năng \\
GitOps  & GitOps & Mô hình triển khai khai báo, đồng bộ trạng thái cụm từ Git \\
IP      & Internet Protocol & Giao thức Internet dùng để định danh và chuyển gói tin \\
K8s     & Kubernetes & Hệ thống điều phối container (deployment, scaling, scheduling) \\
kprobe  & Kernel probe & Điểm chèn (hook) động vào hàm nhân để chạy chương trình eBPF \\
LPM     & Longest Prefix Match & Khớp tiền tố dài nhất (phân loại địa chỉ theo CIDR) \\
MTTR    & Mean Time To Respond & Thời gian trung bình để phản hồi/xử lý sự cố \\
NetworkPolicy & Kubernetes NetworkPolicy & Đối tượng K8s khai báo luật cho/chặn lưu lượng pod (do CNI thực thi) \\
NFR     & Non-Functional Requirement & Yêu cầu phi chức năng \\
NIDS    & Network Intrusion Detection System & Hệ phát hiện xâm nhập ở mức mạng \\
nmap    & Network Mapper & Công cụ quét cổng/dịch vụ mạng phổ biến \\
OCI     & Open Container Initiative & Chuẩn ảnh và runtime container mở \\
RBAC    & Role-Based Access Control & Kiểm soát truy cập theo vai trò \\
RingBuf & Ring Buffer & Bộ đệm vòng của eBPF để chuyển sự kiện lên không gian người dùng \\
SCC     & Strongly Connected Component & Thành phần liên thông mạnh (trong đồ thị có hướng) \\
SIEM    & Security Information and Event Management & Hệ quản lý thông tin và sự kiện an ninh (phân tích nhật ký) \\
SLR     & Systematic Literature Review & Tổng quan tài liệu có hệ thống \\
SOC     & Security Operations Center & Trung tâm vận hành an ninh (giám sát và phản ứng sự cố) \\
SRE     & Site Reliability Engineering & Kỹ nghệ độ tin cậy hệ thống (vận hành) \\
TCP     & Transmission Control Protocol & Giao thức truyền tải hướng kết nối, tin cậy \\
UC      & Use Case & Ca sử dụng (trong phân tích yêu cầu) \\
UML     & Unified Modeling Language & Ngôn ngữ mô hình hóa thống nhất \\
YAML    & YAML Ain't Markup Language & Ngôn ngữ dữ liệu cấu hình dạng phân cấp \\

\end{longtable}

\clearpage

% ============================================================================
% Bang dien giai ky hieu — cac ky hieu dung trong BC2
% ============================================================================

\chapter*{BẢNG DIỄN GIẢI KÝ HIỆU}
\addcontentsline{toc}{chapter}{BẢNG DIỄN GIẢI KÝ HIỆU}

\begin{longtable}{@{}p{3.5cm} p{10.5cm}@{}}

$S$                    & Luồng sự kiện kết nối đầu vào của bộ phát hiện \\
$\tau$                 & Ngưỡng phát hiện quét cổng (số kết nối/cửa sổ; trong đồ án $\tau = 40$) \\
$d,\ w$                & Chiều sâu (số hàm băm) và chiều rộng của Count-Min Sketch ($d=5$, $w=2048$) \\
$M$                    & Ma trận đếm của Count-Min Sketch, kích thước $d \times w$ \\
$h_i$                  & Hàm băm thứ $i$ của Count-Min Sketch, $i = 1, \ldots, d$ \\
$\hat{f}_x$            & Tần suất (số kết nối) ước lượng của nguồn $x$ trong cửa sổ \\
$\varepsilon,\ \delta$ & Tham số sai số và xác suất thất bại của ước lượng Count-Min Sketch \\
$F_0$                  & Số nguồn phân biệt xuất hiện trong cửa sổ (lực lượng tập khóa) \\
$G = (V, E)$           & Đồ thị truyền thông có hướng: đỉnh $V$ là địa chỉ IP, cung $E$ là kết nối \\
$s$                    & Đỉnh nguồn (pod nghi vấn) trong phân tích đồ thị \\
$k$                    & Số bước (hop) tối đa của BFS bán kính ảnh hưởng (trong đồ án $k = 3$) \\
$R_k(s)$               & Bán kính ảnh hưởng: tập pod tới được từ nguồn $s$ trong không quá $k$ bước \\
$p_{50},\ p_{95}$      & Trung vị và phân vị 95 của độ trễ phản hồi \\
$O(\cdot)$             & Ký hiệu Big-O: cận trên về độ phức tạp thời gian/không gian \\

\end{longtable}
